%%%%%%%%%%%%%%%%%%%%%%%%%%%%%%%%%%%%%%%%%%%%%%%%%%%%%%%%%%%%%%%%%%%%%%%%%%%%%%
%  UHV-QuestionBank.tex
%  Universal Human Values (HS248AT) -- consolidated previous-year question
%  bank. Every paper is reproduced in full, in chronological order, with its
%  date, assessment type, duration and marks recorded on its own opening page,
%  and with the per-question marks / CO / BT codes exactly as printed.
%
%  Sources: Improvement CIE 28 Aug 2024; CIE-II 7 May 2025; Improvement CIE
%           Jun 2025 (Mechanical Engineering); SEE Jun/Jul 2025;
%           Improvement CIE 18 Jun 2026.
%
%  Self-contained: standard CTAN packages only, no custom class required.
%  Build:  pdflatex UHV-QuestionBank   (three times, for the TOC and links)
%%%%%%%%%%%%%%%%%%%%%%%%%%%%%%%%%%%%%%%%%%%%%%%%%%%%%%%%%%%%%%%%%%%%%%%%%%%%%%
\documentclass[11pt,a4paper,oneside]{book}

\usepackage[T1]{fontenc}
\usepackage[utf8]{inputenc}
\usepackage{lmodern}
\usepackage{microtype}
\usepackage[left=22mm,right=22mm,top=24mm,bottom=26mm,headsep=8mm]{geometry}
\usepackage[table]{xcolor}
\usepackage{enumitem}
\usepackage{titlesec}
\usepackage{fancyhdr}
\usepackage{booktabs}
\usepackage{longtable}
\usepackage{array}
\usepackage{needspace}
\usepackage{amssymb}
\usepackage[most]{tcolorbox}
\usepackage[hidelinks]{hyperref}

% ------------------------------------------------------------------- colours
\definecolor{ocre}{RGB}{243,102,25}
\definecolor{slate}{RGB}{45,55,72}
\definecolor{slatelight}{RGB}{113,128,150}
\definecolor{rulegrey}{RGB}{214,222,232}
\definecolor{tint}{RGB}{253,246,240}
\definecolor{panel}{RGB}{247,250,252}

\hypersetup{colorlinks=true, linkcolor=slate, urlcolor=ocre,
            pdftitle={Universal Human Values HS248AT -- Previous-Year
                      Question Bank},
            pdfsubject={IV Semester B.E., R V College of Engineering}}

% ---------------------------------------------------------------- structures
\newcolumntype{P}[1]{>{\raggedright\arraybackslash}p{#1}}
\newcolumntype{C}[1]{>{\centering\arraybackslash}p{#1}}

\setlength{\tabcolsep}{4pt}
\setlength{\arrayrulewidth}{0.4pt}

% ------------------------------------------------------------------ headings
\newcommand{\chapkind}{PAPER}

\titleformat{\chapter}[display]
  {\sffamily\bfseries\color{slate}}
  {\raggedright\normalsize\sffamily\color{ocre}\chapkind\ \thechapter}
  {6pt}
  {\LARGE\raggedright}
  [\vspace{4pt}{\color{ocre}\rule{\textwidth}{1.6pt}}]
\titlespacing*{\chapter}{0pt}{-24pt}{18pt}

\newcommand{\qbsecbox}[1]{%
  \colorbox{slate}{\parbox{\dimexpr\textwidth-2\fboxsep\relax}%
    {\sffamily\bfseries\large\color{white}\strut #1\strut}}}
\titleformat{\section}[block]{}{}{0pt}{\qbsecbox}
\titlespacing*{\section}{0pt}{16pt}{8pt}

\titleformat{\subsection}[block]
  {\sffamily\bfseries\color{slate}\large}{}{0pt}{}
  [\vspace{-7pt}{\color{ocre}\rule{0.16\textwidth}{1pt}}]
\titlespacing*{\subsection}{0pt}{14pt}{4pt}

\pagestyle{fancy}
\fancyhf{}
\renewcommand{\headrulewidth}{0.4pt}
\fancyhead[L]{\footnotesize\sffamily\color{slatelight}HS248AT \textbullet\
  Universal Human Values}
\fancyhead[R]{\footnotesize\sffamily\color{slatelight}\leftmark}
\fancyfoot[C]{\footnotesize\sffamily\color{slatelight}\thepage}
\renewcommand{\chaptermark}[1]{\markboth{#1}{}}
\fancypagestyle{plain}{\fancyhf{}%
  \fancyfoot[C]{\footnotesize\sffamily\color{slatelight}\thepage}%
  \renewcommand{\headrulewidth}{0pt}}

% --------------------------------------------------------------------- macros
% a fill-in-the-blank rule, as printed on the paper
\newcommand{\blank}{\underline{\hspace{2.6em}}}

% inline options, for short choices; \quad is breakable, so long rows wrap
\newcommand{\optsinline}[4]{\par\vspace{3pt}%
  {\small\color{slate}(a)~#1\quad (b)~#2\quad (c)~#3\quad (d)~#4}}

% stacked options, for choices that are clauses rather than words
\newenvironment{optlist}
  {\par\vspace{2pt}\begingroup\small\color{slate}%
   \begin{enumerate}[label={\color{slatelight}(\alph*)},leftmargin=2.0em,
     itemsep=1pt,topsep=2pt,parsep=0pt]}
  {\end{enumerate}\endgroup\vspace{1pt}}

% ------------------------------------------------------------ the question grid
\newcommand{\hdr}[1]{{\sffamily\bfseries\color{white}\strut #1}}

\newenvironment{qtable}
  {\renewcommand{\arraystretch}{1.28}%
   \arrayrulecolor{rulegrey}%
   \setlength{\LTpre}{6pt}\setlength{\LTpost}{4pt}%
   \begin{longtable}{C{11mm} P{108mm} C{10mm} C{12mm} C{9mm}}
   \rowcolor{slate}
   \hdr{Q.\,No} & \hdr{Question} & \hdr{M} & \hdr{CO} & \hdr{BT}
   \\ \hline \endfirsthead
   \rowcolor{slate}
   \hdr{Q.\,No} & \hdr{Question} & \hdr{M} & \hdr{CO} & \hdr{BT}
   \\ \hline \endhead}
  {\end{longtable}}

% the theme index grid used in Appendix A
\newenvironment{idxtable}
  {\renewcommand{\arraystretch}{1.3}%
   \arrayrulecolor{rulegrey}%
   \setlength{\LTpre}{6pt}\setlength{\LTpost}{4pt}%
   \begin{longtable}{P{104mm} P{40mm} C{13mm}}
   \rowcolor{slate}
   \hdr{Question} & \hdr{Paper} & \hdr{Marks} \\ \hline \endfirsthead
   \rowcolor{slate}
   \hdr{Question} & \hdr{Paper} & \hdr{Marks} \\ \hline \endhead}
  {\end{longtable}}

% question number cell
\newcommand{\qn}[1]{{\sffamily\bfseries\color{ocre}#1}}
% the "OR" band inside a Part B grid
\newcommand{\orband}{\multicolumn{5}{c}{\cellcolor{tint}%
  \sffamily\bfseries\color{ocre}\strut OR} \\ \hline}

% ------------------------------------------------------- paper metadata panel
\newenvironment{papermeta}
  {\begin{tcolorbox}[colback=panel,colframe=rulegrey,boxrule=0.8pt,
     sharp corners,left=8pt,right=8pt,top=7pt,bottom=7pt,
     before skip=4pt,after skip=10pt]
   \renewcommand{\arraystretch}{1.25}%
   \begin{tabular}{@{}P{24mm}P{48mm}@{\hspace{6mm}}P{24mm}P{46mm}@{}}}
  {\end{tabular}\end{tcolorbox}}
\newcommand{\mk}[1]{{\footnotesize\sffamily\bfseries\color{slatelight}%
  \MakeUppercase{#1}}}
\newcommand{\mv}[1]{{\color{slate}#1}}

% the marks-distribution footer table, as printed on the CIE papers
\newenvironment{mdtable}
  {\renewcommand{\arraystretch}{1.3}%
   \arrayrulecolor{rulegrey}%
   \begin{tabular}{@{}P{38mm} C{11mm}C{11mm}C{11mm}C{11mm}
                        C{9mm}C{9mm}C{9mm}C{9mm}C{9mm}C{9mm}@{}}
   \rowcolor{slate}
   \hdr{Particulars} & \hdr{CO1} & \hdr{CO2} & \hdr{CO3} & \hdr{CO4} &
   \hdr{L1} & \hdr{L2} & \hdr{L3} & \hdr{L4} & \hdr{L5} & \hdr{L6}
   \\ \hline}
  {\end{tabular}}

\newenvironment{note}
  {\begin{tcolorbox}[colback=tint,colframe=ocre,boxrule=0.8pt,sharp corners,
     fontupper=\small,breakable,before skip=8pt,after skip=8pt]}
  {\end{tcolorbox}}

\newenvironment{abul}
  {\begin{itemize}[leftmargin=1.5em,itemsep=2pt,topsep=3pt,parsep=0pt,
                   label=\textcolor{ocre}{\scriptsize$\blacksquare$}]}
  {\end{itemize}}

\setlength{\parindent}{0pt}
\setlength{\parskip}{3pt}
\setlength{\emergencystretch}{3em}

%%%%%%%%%%%%%%%%%%%%%%%%%%%%%%%%%%%%%%%%%%%%%%%%%%%%%%%%%%%%%%%%%%%%%%%%%%%%%%
\begin{document}
\frontmatter
\pagestyle{empty}

% ---------------------------------------------------------------- title page
\begin{center}
\vspace*{32mm}
{\sffamily\color{ocre}\large HS248AT \textbullet\ IV SEMESTER B.E.}\\[3mm]
{\color{rulegrey}\rule{0.55\textwidth}{1pt}}\\[9mm]
{\sffamily\bfseries\color{slate}\Huge UNIVERSAL HUMAN VALUES}\\[7mm]
{\sffamily\bfseries\color{slate}\LARGE Previous-Year Question Bank}\\[3mm]
{\sffamily\color{slatelight}\Large Every paper reproduced in full, dated and
 classified}\\[12mm]
{\color{rulegrey}\rule{0.55\textwidth}{1pt}}\\[9mm]
{\sffamily\color{slate}\large Five papers \textbullet\ 59 questions
 \textbullet\ August 2024 to June 2026}\\[3mm]
{\sffamily\color{slatelight}\large Semester End Examination \quad\textbullet\quad
 CIE-II \quad\textbullet\quad Improvement CIE}\\[16mm]
{\sffamily\color{slatelight} Department of Industrial Engineering and
 Management}\\[1mm]
{\sffamily\color{slatelight} R V College of Engineering, Bengaluru}\\
\vfill
{\footnotesize\sffamily\color{slatelight}
 Unofficial student-compiled study material. Not issued, endorsed or verified
 by the Department or by RV College of Engineering. Every question is
 reproduced from the question paper as printed.}
\vspace*{12mm}
\end{center}
\clearpage

% -------------------------------------------------------------- how arranged
\chapter*{How This Book Is Arranged}
\addcontentsline{toc}{chapter}{How This Book Is Arranged}
\markboth{How This Book Is Arranged}{}
\thispagestyle{plain}

This book collects every previous-year question recorded for HS248AT ---
\textbf{59 questions} in total, \textbf{30 objective or short-answer} and
\textbf{29 descriptive} --- drawn from five papers set between August 2024 and
June 2026.

\vspace{2pt}
The questions are printed \textbf{paper by paper, in chronological order},
not merged under topic headings. A question paper is a document with a shape:
how many marks Part~A carries, how many descriptive questions must be
attempted, which choices are offered. That shape is itself worth revising, and
it is lost the moment the questions are shuffled. Each chapter therefore opens
with the paper's own header --- date, assessment type, duration, maximum
marks and, where the paper printed them, the instructions to candidates and
the marks-distribution table --- followed by the questions in their original
order and numbering.

\vspace{2pt}
Appendix~A puts back what the chronological arrangement leaves out. It groups
all 59 questions by theme, so everything ever asked on the four orders, or on
trust and respect, can be read together. Appendix~B lists the questions that
have been set more than once; those repeats are the single best use of
revision time.

\subsection*{The papers in this book}
\addcontentsline{toc}{section}{The papers in this book}

\renewcommand{\arraystretch}{1.35}
\arrayrulecolor{rulegrey}
\begin{tabular}{@{}C{9mm} P{44mm} P{30mm} C{18mm} C{20mm} C{17mm}@{}}
\rowcolor{slate}
\hdr{\#} & \hdr{Assessment} & \hdr{Date} & \hdr{Max M} & \hdr{Time} &
\hdr{Questions} \\ \hline
\qn{1} & Improvement CIE --- Quiz and Test & 28 August 2024 & 25 & 60 min &
  10 \\ \hline
\qn{2} & CIE-II --- Quiz and Test & 7 May 2025 & 25 $+$ 5 & 60 min &
  10 \\ \hline
\qn{3} & Improvement CIE --- Quiz and Test & June 2025 & 30 & 60 min &
  9 \\ \hline
\qn{4} & \textbf{Semester End Examination} --- Regular / Supplementary &
  June/July 2025 & 50 & 2 hours & 20 \\ \hline
\qn{5} & Improvement CIE --- Quiz and Test & 18 June 2026 & 05 $+$ 25 &
  10 $+$ 60 min & 10 \\ \hline
\end{tabular}

\subsection*{Reading the three code columns}
\addcontentsline{toc}{section}{Reading the three code columns}

Every question in this book carries the three codes its paper printed beside
it, in the same three columns:

\begin{abul}
\item \textbf{M} --- the marks the question carries.
\item \textbf{CO} --- the Course Outcome the question is mapped to. The papers
      print the code only, CO1 to CO4; the full statement of each outcome is
      given in the course handout and is not reproduced here.
\item \textbf{BT} --- the Bloom's Taxonomy level, on the standard six-level
      scale: \textbf{L1} remember, \textbf{L2} understand, \textbf{L3} apply,
      \textbf{L4} analyse, \textbf{L5} evaluate, \textbf{L6} create. No paper
      in this collection sets anything above L4.
\end{abul}

The BT code is the most useful of the three when revising. An L1 question wants
a definition reproduced; an L2 question wants it explained in your own words;
an L3 question wants it used on a case; an L4 question wants it taken apart and
compared. The same topic --- conformance in the four orders, say --- has been
set at L4 on one paper and at L2 on another, and the two answers are not the
same length or the same shape.

\begin{note}
\textbf{\sffamily A note on faithfulness.} Questions are transcribed as
printed, including the spelling each paper used where it varies between them
(\emph{swabhava} in August 2024, \emph{svabhava} in June/July 2025) and
including the occasional grammatical slip in the original. Nothing has been
silently corrected, abridged or reworded, and no question has been dropped.
The same applies to the papers' own arithmetic: the August 2024 paper sets 30
marks under a header that says 25, and that is reproduced as printed, with the
discrepancy noted where it occurs rather than tidied away.
\end{note}

\clearpage
\tableofcontents

\mainmatter
\pagestyle{fancy}

%%%%%%%%%%%%%%%%%%%%%%%%%%%%%%%%%%%%%%%%%%%%%%%%%%%%%% PAPER 1 -- 28 Aug 2024
\chapter{Improvement CIE --- 28 August 2024}

\begin{papermeta}
\mk{Course code}   & \mv{HS248AT}                    &
\mk{Maximum marks} & \mv{25}                          \\
\mk{Semester}      & \mv{IV}                          &
\mk{Duration}      & \mv{60 minutes}                  \\
\mk{Academic year} & \mv{2023--2024, Even Semester}   &
\mk{Assessment}    & \mv{Improvement CIE --- Quiz and Test} \\
\mk{Date}          & \mv{28 August 2024}              &
\mk{Structure}     & \mv{Part A 5 M \textbf{+} Part B 25 M} \\
\end{papermeta}

\section[Part A --- Objective Questions]{Part A --- Objective Questions
  \hfill 5 Marks}

\begin{qtable}
\qn{1} & What is the fundamental principle of harmony in nature and existence?
  \optsinline{Competition}{Coexistence}{Dominance}{Isolation}
  & 1 & CO3 & L1 \\ \hline
\qn{2} & What is the role of human beings in maintaining harmony in nature?
  \begin{optlist}
  \item Exploiting natural resources
  \item Preserving and nurturing nature
  \item Dominating other species
  \item Isolating from the environment
  \end{optlist}
  & 1 & CO2 & L2 \\ \hline
\qn{3} & Which of the following best describes the Physical Order?
  \begin{optlist}
  \item Includes only human-made objects
  \item Comprises natural elements like soil, water, air, and minerals
  \item Refers exclusively to living beings
  \item Focuses on mental and spiritual aspects
  \end{optlist}
  & 1 & CO2 & L1 \\ \hline
\qn{4} & What distinguishes the Human Order from the other three orders?
  \begin{optlist}
  \item Physical strength
  \item Intelligence and moral values
  \item Ability to photosynthesize
  \item Instinctual behaviour only
  \end{optlist}
  & 1 & CO4 & L1 \\ \hline
\qn{5} & In the context of coexistence, what is the meaning of `prosperity'?
  \begin{optlist}
  \item Material wealth accumulation
  \item Holistic well-being and happiness
  \item Physical strength and power
  \item Technological superiority
  \end{optlist}
  & 1 & CO3 & L1 \\ \hline
\end{qtable}

\section[Part B --- Descriptive Questions]{Part B --- Descriptive Questions
  \hfill 25 Marks}

\begin{qtable}
\qn{1} & Existence is co-existence of mutually interacting units in
  all-pervasive space. Explain & 5 & CO3 & L2 \\ \hline
\qn{2} & What is the natural characteristics (swabhava) of human order?
  Explain & 5 & CO2 & L1 \\ \hline
\qn{3} & What are the four orders of nature? Briefly explain them
  & 5 & CO4 & L2 \\ \hline
\qn{4} & How can we say that `nature is Self-Organized and in space
  Self-Organization Is Available.' & 5 & CO4 & L2 \\ \hline
\qn{5} & Explain the difference and similarities between Pranic order and
  animal order. What is the relation between the two orders? & 5 & CO2 & L1
  \\ \hline
\end{qtable}

\subsection*{Marks distribution as printed}

\begin{mdtable}
Test --- Max Marks & xx & 12 & 07 & 11 & 14 & 16 & xx & xx & xx & xx \\ \hline
\end{mdtable}

\vspace{4pt}
Four cells were left as the template's own \texttt{xx} placeholder and are
reproduced that way. They are all cells that would have been zero: no question
on this paper is mapped to CO1, and none is set above L2.

\begin{note}
\textbf{\sffamily The paper does not add up to its own maximum.} The header
states a maximum of \textbf{25 marks}, but Part~A carries 5 and the five
Part~B questions carry 5 each --- \textbf{30 in total}. The footer table above
agrees with the questions, not with the header: its CO row totals 30
($12+7+11$) and its BT row totals 30 ($14+16$), and both reconcile exactly
with the per-question codes in the two grids. The discrepancy is in the
original and has been left as printed.
\end{note}

%%%%%%%%%%%%%%%%%%%%%%%%%%%%%%%%%%%%%%%%%%%%%%%%%%%%%%% PAPER 2 -- 7 May 2025
\chapter{CIE-II --- 7 May 2025}

\begin{papermeta}
\mk{Course code}   & \mv{HS248AT}                    &
\mk{Maximum marks} & \mv{25 $+$ 5}                    \\
\mk{Semester}      & \mv{IV}                          &
\mk{Duration}      & \mv{60 minutes}                  \\
\mk{Academic year} & \mv{2024--2025, Even Semester}   &
\mk{Assessment}    & \mv{CIE-II --- Quiz and Test}    \\
\mk{Date}          & \mv{7 May 2025}                  &
\mk{Structure}     & \mv{Part A 5 M \textbf{+} Part B 25 M} \\
\end{papermeta}

This is the only regular CIE in the collection --- the other four papers are
improvement or semester-end sittings. Its subject matter is correspondingly
narrower: every question on it comes from harmony in the family and the values
in relationship, which is the block of the syllabus a second CIE covers.

\section[Part A --- Objective Questions]{Part A --- Objective Questions
  \hfill 5 Marks}

\begin{qtable}
\qn{1} & Which of the following is considered the foundation of a strong human
  relationship?
  \optsinline{Intelligence}{Trust}{Wealth}{Appearance}
  & 1 & CO1 & L1 \\ \hline
\qn{2} & Disrespect Arising out of
  \begin{optlist}
  \item Lack of communication
  \item Lack of knowledge
  \item Differentiation leading to Discrimination
  \item Lack of Wealth
  \end{optlist}
  & 1 & CO2 & L1 \\ \hline
\qn{3} & Care is the feeling
  \begin{optlist}
  \item of responsibility and commitment for nurturing and protection of the
        Body of my relative.
  \item of responsibility and providing facility for nurturing and protection
        of the Body of my relative
  \item of responsibility of making the other feel comfortable with respect to
        Body of my relative
  \item of showing affection to the Body of my relative
  \end{optlist}
  & 1 & CO2 & L1 \\ \hline
\qn{4} & Fulfilment of relationship means
  \begin{optlist}
  \item Making effort for self development, i.e. development of one's own
        competence
  \item Making effort for mutual development, i.e. development of one's own
        competence and being of help to the other in developing their
        competence
  \item Making effort for mutual understanding, i.e. understanding of one's
        own competence and being of help to the other in understanding their
        competence
  \item Making effort for mutual in acquiring prosperity
  \end{optlist}
  & 1 & CO3 & L1 \\ \hline
\qn{5} & With the complete understanding of respect
  \begin{optlist}
  \item we can see for every individual on the earth that we all are the same
        in terms of intelligent, prosperity and status.
  \item we can see for every individual on the earth that we all are the same
        in terms of Knowledge, family, society.
  \item we can see for every individual on the earth that we all are the same
        in terms of competence, prosperity and status.
  \item we can see for every individual on the earth that we all are the same
        in terms of intention, program and potential.
  \end{optlist}
  & 1 & CO1 & L1 \\ \hline
\end{qtable}

\subsection*{Part A --- marks distribution as printed}

\begin{mdtable}
Test --- Max Marks & 2 & 2 & 1 & --- & 5 & --- & --- & --- & --- & --- \\
\hline
\end{mdtable}

\section[Part B --- Descriptive Questions]{Part B --- Descriptive Questions
  \hfill 25 Marks}

\begin{qtable}
\qn{1} & Justify the following (i)~Relation ship is between one Self (I1) and
  another Self (I2), (ii)~feelings in relationship --- in one Self (I1) for the
  other Self (I2) & 2$+$3 & CO1 & L1 \\ \hline
\qn{2} & Distinguish between Intention and Competence with respect Trust with
  suitable example & 5 & CO3 & L2 \\ \hline
\qn{3} & The complete content of respect is to be able to see that `the other
  is similar to me and we are complementary'. Define complementary in this
  context. & 5 & CO1 & L1 \\ \hline
\qn{4} & Disrespect arises out of differentiation leading to discrimination on
  the basis of body, physical facility or beliefs. illustrate With an example
  & 5 & CO2 & L2 \\ \hline
\qn{5} & Illustrate the natural acceptance values in relationship with an
  example for each (i)~Reverence (ii)~Glory & 5 & CO3 & L1 \\ \hline
\end{qtable}

\subsection*{Part B --- marks distribution as printed}

\begin{mdtable}
Test --- Max Marks & 10 & 10 & 5 & --- & 15 & 10 & --- & --- & --- & --- \\
\hline
\end{mdtable}

%%%%%%%%%%%%%%%%%%%%%%%%%%%%%%%%%%%%%%%%%%%%%%%%%%%%%%% PAPER 3 -- June 2025
\chapter{Improvement CIE --- June 2025}

\begin{papermeta}
\mk{Course code}   & \mv{HS248AT}                    &
\mk{Maximum marks} & \mv{30}                          \\
\mk{Semester}      & \mv{IV}                          &
\mk{Duration}      & \mv{60 minutes}                  \\
\mk{Academic year} & \mv{2024--2025, Even Semester}   &
\mk{Assessment}    & \mv{Improvement CIE --- Quiz and Test} \\
\mk{Date}          & \mv{June 2025}                   &
\mk{Structure}     & \mv{Quiz 5 M \textbf{+} Test 25 M}  \\
\mk{Set by}        & \mv{Department of Mechanical Engineering} &
                   & \\
\end{papermeta}

\section[Quiz]{Quiz \hfill 5 Marks}

\begin{qtable}
\qn{1} & The participation of the human being in ensuring the role of physical
  facility in nurture, protection and providing means for the body is called
  its \blank
  \optsinline{Utility value}{Artistic value}{Activity}{Energy}
  & 1 & CO1 & L2 \\ \hline
\qn{2} & When something is active or has activity, it is called as \blank
  \optsinline{space}{unit}{energy}{value}
  & 1 & CO2 & L3 \\ \hline
\qn{3} & The following are the characteristics of space except \blank
  \optsinline{Unlimited}{No activity}{Energy in equilibrium}{Self-organized}
  & 1 & CO1 & L1 \\ \hline
\qn{4} & The natural characteristic of material order is \blank
  \optsinline{Composition}{Decomposition}{Composition/Decomposition}%
  {Nurture/worsen}
  & 1 & CO1 & L1 \\ \hline
\qn{5} & Conformance of material order is named as \blank
  \optsinline{constitution conformance}{Seed conformance}{Breed conformance}%
  {sanskaar conformance}
  & 1 & CO3 & L3 \\ \hline
\end{qtable}

\section[Test]{Test \hfill 25 Marks}

\begin{qtable}
\qn{1} & Explain the classification of units into Four Orders and the harmony
  among them with the help of a schematic diagram. & 6 & CO1 & L2 \\ \hline
\qn{2} & What do you mean by `conformance'? Explain the conformance in the
  four orders. & 6 & CO2 & L4 \\ \hline
\qn{3} & Explain how there is recyclability and self-regulation in nature.
  & 6 & CO3 & L3 \\ \hline
\qn{4} & Elaborate the terms `units' and `space' and mention their
  characteristics. & 7 & CO4 & L4 \\ \hline
\end{qtable}

\subsection*{Marks distribution as printed}

\begin{mdtable}
Test --- Marks & 9 & 7 & 7 & 7 & 3 & 7 & 7 & 13 & --- & --- \\ \hline
\end{mdtable}

\vspace{4pt}
Both rows total 30, which is the paper's maximum. Note that L4 alone carries
13 of the 30 marks: this is an analytical paper, and the six- and seven-mark
questions on conformance and on units and space are where those marks sit.

%%%%%%%%%%%%%%%%%%%%%%%%%%%%%%%%%%%%%%%%%%%%%%%%%%% PAPER 4 -- SEE Jun/Jul 2025
\chapter{Semester End Examination --- June/July 2025}
\chaptermark{SEE --- June/July 2025}

\begin{papermeta}
\mk{Course code}   & \mv{HS248AT}                    &
\mk{Maximum marks} & \mv{50}                          \\
\mk{Semester}      & \mv{IV}                          &
\mk{Duration}      & \mv{2 hours}                     \\
\mk{Examination}   & \mv{IV Semester B.E. Regular / Supplementary} &
\mk{Assessment}    & \mv{Semester End Examination}    \\
\mk{Date}          & \mv{June/July 2025}              &
\mk{Applies to}    & \mv{Common to all programs}      \\
\end{papermeta}

\begin{note}
\textbf{\sffamily Instructions to the students, as printed.}
\begin{enumerate}[leftmargin=1.6em,itemsep=2pt,topsep=4pt]
\item Answer all questions from Part A. Part A questions should be answered in
      first three pages of the answer book only.
\item Answer THREE full questions from Part B. In Part B question number 2 is
      compulsory. Answer any one full question from 3 and 4 \& 5 and 6.
\end{enumerate}
\end{note}

The blueprint is worth reading off the instructions before the questions.
Part~A is ten compulsory one-mark objective questions. Part~B sets
\textbf{Q2 compulsory} (14 marks), then \textbf{one of Q3 or Q4} (13 marks) and
\textbf{one of Q5 or Q6} (13 marks) --- three full questions, 40 marks, and two
genuine choices to make in the hall.

\section[Part A --- Compulsory Objective Questions]{Part A --- Compulsory
  Objective Questions \hfill 10 Marks}

\begin{qtable}
\qn{1.1} & How does ensuring justice in relationships contribute to society?
  \begin{optlist}
  \item It leads to fearlessness and co-existence
  \item It leads to prosperity and right understanding
  \item It promotes self-regulation and health
  \item It ensures better exchange and storage
  \end{optlist}
  & 01 & CO3 & L1 \\ \hline
\qn{1.2} & Imagination is continuous with
  \optsinline{Time}{Soul}{Body}{None}
  & 01 & CO1 & L1 \\ \hline
\qn{1.3} & Identify the solution which helps human being to transform from
  animal consciousness to human consciousness.
  \optsinline{Right understanding}{Realization}{Value education}%
  {Physical facilities.}
  & 01 & CO1 & L1 \\ \hline
\qn{1.4} & Which of the following is a fundamental human aspiration related to
  relationships?
  \begin{optlist}
  \item Physical comfort.
  \item Financial security.
  \item Mutual understanding and fulfillment.
  \item Accumulating wealth.
  \end{optlist}
  & 01 & CO2 & L2 \\ \hline
\qn{1.5} & Which of the following capacity leads to desires
  \optsinline{Power}{Expectation}{Realization}{Thoughts}
  & 01 & CO2 & L1 \\ \hline
\qn{1.6} & Which of the following is considered the foundation of a strong
  human relationship?
  \optsinline{Intelligence}{Trust}{Wealth}{Appearance}
  & 01 & CO1 & L2 \\ \hline
\qn{1.7} & Respect in a relationship means:
  \begin{optlist}
  \item Agreeing to everything
  \item Controlling the other person's actions
  \item Valuing the other person's opinions and boundaries
  \item Always trying to win argument.
  \end{optlist}
  & 01 & CO3 & L1 \\ \hline
\qn{1.8} & The participation of the human being in ensuring the role of
  physical facility to help and preserve its utility is called its \blank
  \optsinline{Utility value}{Artistic value}{Activity}{Energy}
  & 01 & CO3 & L2 \\ \hline
\qn{1.9} & A harmonious world is created by values at 4 levels. These are:
  \begin{optlist}
  \item Home, family, society, country
  \item Individual, family, society, universe
  \item School, home, office, temple
  \item None of these
  \end{optlist}
  & 01 & CO1 & L1 \\ \hline
\qn{1.10} & When nature is submerged in space it is known as
  \optsinline{Conformance}{Acceptance}{Mixing}{Co-existence}
  & 01 & CO3 & L2 \\ \hline
\end{qtable}

\section[Part B --- Descriptive Questions]{Part B --- Descriptive Questions
  \hfill 40 Marks}

\subsection*{Question 2 --- compulsory}

\begin{qtable}
\qn{2a} & What do you understand by the value of an entity? What is the value
  of a human being? & 08 & CO1 & L2 \\ \hline
\qn{2b} & For human being physical facility is necessary, but relationship is
  also necessary: Justify with an example comparing with animal order.
  & 06 & CO2 & L4 \\ \hline
\end{qtable}

\subsection*{Questions 3 and 4 --- answer any one full question}

\begin{qtable}
\qn{3a} & ``Discuss how right understanding in relationships contributes to
  harmony in the family. Analyze the role of trust, respect, and
  responsibility in maintaining family unity, and provide examples to support
  your answer.'' & 07 & CO4 & L4 \\ \hline
\qn{3b} & Articulate the four goals of human beings living in a society.
  Critically examine the state of the society today in context with the
  fulfilment of a comprehensive human goal. & 06 & CO3 & L3 \\ \hline
\orband
\qn{4a} & How do justice and preservation contribute to a harmonious society?
  Analyze the shortcomings of focusing solely on punishment in cases of
  injustice, and explain how true preservation ensures prosperity and
  coexistence with nature. & 06 & CO3 & L4 \\ \hline
\qn{4b} & Explain the concept of harmony in the family. Why is it important,
  and how can it be achieved through right understanding and mutual
  relationships? & 07 & CO3 & L2 \\ \hline
\end{qtable}

\subsection*{Questions 5 and 6 --- answer any one full question}

\begin{qtable}
\qn{5a} & What is sanskaar? Explain its effects or the conformance of the
  human order. & 06 & CO4 & L4 \\ \hline
\qn{5b} & What is the svabhava (natural characteristic) of a unit? Mention the
  natural characteristics of the material, pranic, animal and human orders.
  Give examples & 07 & CO4 & L3 \\ \hline
\orband
\qn{6a} & What do you mean by `conformance'? Explain the conformance in the
  four orders. & 07 & CO3 & L4 \\ \hline
\qn{6b} & Explain the meaning of co-existence. & 06 & CO3 & L3 \\ \hline
\end{qtable}

%%%%%%%%%%%%%%%%%%%%%%%%%%%%%%%%%%%%%%%%%%%%%%%%%%%%%% PAPER 5 -- 18 Jun 2026
\chapter{Improvement CIE --- 18 June 2026}

\begin{papermeta}
\mk{Course code}   & \mv{HS248AT}                    &
\mk{Maximum marks} & \mv{05 $+$ 25}                   \\
\mk{Semester}      & \mv{IV}                          &
\mk{Duration}      & \mv{10 min $+$ 60 min}           \\
\mk{Date}          & \mv{18 June 2026}                &
\mk{Assessment}    & \mv{Improvement CIE --- Quiz $+$ Test} \\
\mk{Structure}     & \mv{Part A 5 M \textbf{+} Part B 25 M} &
\mk{Format}        & \mv{Short answer, not multiple choice} \\
\end{papermeta}

This paper breaks the pattern of the earlier improvement CIEs in one respect
worth noting: \textbf{Part~A is not multiple choice}. It sets five one-mark
short-answer questions to be written in ten minutes --- definitions, lists and
one-line differentiations. There is nothing to guess at.

\section[Part A --- Short Answer Questions]{Part A --- Short Answer Questions
  \hfill 5 Marks}

\begin{qtable}
\qn{1} & Give any two definitions for the term co-existence.
  & 1 & CO1 & L1 \\ \hline
\qn{2} & State the four orders of nature. & 1 & CO2 & L1 \\ \hline
\qn{3} & Identify the svabhava\,/\,value of the self (`I') in human beings.
  & 1 & CO2 & L2 \\ \hline
\qn{4} & Summarize 2 things that human beings have to do to realize the mutual
  fulfilment in nature. & 1 & CO3 & L2 \\ \hline
\qn{5} & Differentiate between Units and Space. & 1 & CO4 & L2 \\ \hline
\end{qtable}

\section[Part B --- Descriptive Questions]{Part B --- Descriptive Questions
  \hfill 25 Marks}

\begin{qtable}
\qn{1} & Describe how harmony is achieved among the four orders of nature.
  & 5 & CO4 & L2 \\ \hline
\qn{2} & Justify that requirement of any order is already available in
  abundance in nature. & 5 & CO3 & L4 \\ \hline
\qn{3} & Demonstrate how interconnectedness, self-regulation and mutual
  fulfilment are among the four orders of nature with an example.
  & 5 & CO2 & L3 \\ \hline
\qn{4} & Summarize how the domain of consciousness helps in accomplishing
  development or permanent change. & 5 & CO4 & L2 \\ \hline
\qn{5} & Outline 2 types of units with examples. & 5 & CO3 & L2 \\ \hline
\end{qtable}

\subsection*{Marks distribution as printed}

\begin{mdtable}
Test --- Max Marks & 1 & 7 & 11 & 11 & 2 & 18 & 5 & 5 & --- & --- \\ \hline
\end{mdtable}

\vspace{4pt}
Both rows total 30. L2 carries 18 of the 30 marks, which is the mirror image of
the June 2025 paper: this one rewards clear explanation rather than analysis.

%%%%%%%%%%%%%%%%%%%%%%%%%%%%%%%%%%%%%%%%%%%%%%%%%%%%%%%%%%%%%%%%%%% APPENDICES
\appendix
\renewcommand{\chapkind}{APPENDIX}

\chapter{Question Index by Theme}

The chronological arrangement of this book keeps each paper whole. This
appendix inverts it: all 59 questions are listed again under the topic they
examine, so that everything ever asked on one theme can be revised together.
Each entry gives the paper, the question number and the marks.

Five themes account for the whole collection, and their weights are uneven.
The unevenness is the point --- \emph{the four orders} and \emph{harmony in
nature and existence} together carry more than half of every mark ever set on
this course.

\subsection*{A.1\enspace The four orders, units and space}
\addcontentsline{toc}{section}{A.1 The four orders, units and space}

\begin{idxtable}
Which of the following best describes the Physical Order? &
  Improvement CIE, 28 Aug 2024, Part A Q3 & 1 \\ \hline
What distinguishes the Human Order from the other three orders? &
  Improvement CIE, 28 Aug 2024, Part A Q4 & 1 \\ \hline
What is the natural characteristics (swabhava) of human order? &
  Improvement CIE, 28 Aug 2024, Part B Q2 & 5 \\ \hline
What are the four orders of nature? Briefly explain them &
  Improvement CIE, 28 Aug 2024, Part B Q3 & 5 \\ \hline
Difference and similarities between Pranic order and animal order &
  Improvement CIE, 28 Aug 2024, Part B Q5 & 5 \\ \hline
When something is active or has activity, it is called as \blank &
  Improvement CIE, Jun 2025, Quiz Q2 & 1 \\ \hline
The following are the characteristics of space except \blank &
  Improvement CIE, Jun 2025, Quiz Q3 & 1 \\ \hline
The natural characteristic of material order is \blank &
  Improvement CIE, Jun 2025, Quiz Q4 & 1 \\ \hline
Conformance of material order is named as \blank &
  Improvement CIE, Jun 2025, Quiz Q5 & 1 \\ \hline
Classification of units into Four Orders and the harmony among them &
  Improvement CIE, Jun 2025, Test Q1 & 6 \\ \hline
What do you mean by `conformance'? Explain it in the four orders &
  Improvement CIE, Jun 2025, Test Q2 & 6 \\ \hline
Elaborate the terms `units' and `space' and their characteristics &
  Improvement CIE, Jun 2025, Test Q4 & 7 \\ \hline
What is sanskaar? Explain the conformance of the human order &
  SEE, Jun/Jul 2025, Q5a & 06 \\ \hline
What is the svabhava of a unit? Natural characteristics of the four orders &
  SEE, Jun/Jul 2025, Q5b & 07 \\ \hline
What do you mean by `conformance'? Explain it in the four orders &
  SEE, Jun/Jul 2025, Q6a & 07 \\ \hline
State the four orders of nature. &
  Improvement CIE, 18 Jun 2026, Part A Q2 & 1 \\ \hline
Identify the svabhava\,/\,value of the self (`I') in human beings. &
  Improvement CIE, 18 Jun 2026, Part A Q3 & 1 \\ \hline
Differentiate between Units and Space. &
  Improvement CIE, 18 Jun 2026, Part A Q5 & 1 \\ \hline
Describe how harmony is achieved among the four orders of nature. &
  Improvement CIE, 18 Jun 2026, Part B Q1 & 5 \\ \hline
Interconnectedness, self-regulation and mutual fulfilment among the orders &
  Improvement CIE, 18 Jun 2026, Part B Q3 & 5 \\ \hline
Outline 2 types of units with examples. &
  Improvement CIE, 18 Jun 2026, Part B Q5 & 5 \\ \hline
\end{idxtable}

\subsection*{A.2\enspace Harmony in nature, existence and co-existence}
\addcontentsline{toc}{section}{A.2 Harmony in nature, existence and
  co-existence}

\begin{idxtable}
Fundamental principle of harmony in nature and existence &
  Improvement CIE, 28 Aug 2024, Part A Q1 & 1 \\ \hline
Role of human beings in maintaining harmony in nature &
  Improvement CIE, 28 Aug 2024, Part A Q2 & 1 \\ \hline
In the context of coexistence, what is the meaning of `prosperity'? &
  Improvement CIE, 28 Aug 2024, Part A Q5 & 1 \\ \hline
Existence is co-existence of mutually interacting units in all-pervasive
  space & Improvement CIE, 28 Aug 2024, Part B Q1 & 5 \\ \hline
`Nature is Self-Organized and in space Self-Organization Is Available' &
  Improvement CIE, 28 Aug 2024, Part B Q4 & 5 \\ \hline
Explain how there is recyclability and self-regulation in nature. &
  Improvement CIE, Jun 2025, Test Q3 & 6 \\ \hline
When nature is submerged in space it is known as &
  SEE, Jun/Jul 2025, Q1.10 & 01 \\ \hline
Explain the meaning of co-existence. & SEE, Jun/Jul 2025, Q6b & 06 \\ \hline
Give any two definitions for the term co-existence. &
  Improvement CIE, 18 Jun 2026, Part A Q1 & 1 \\ \hline
Two things human beings must do to realize mutual fulfilment in nature &
  Improvement CIE, 18 Jun 2026, Part A Q4 & 1 \\ \hline
Requirement of any order is already available in abundance in nature &
  Improvement CIE, 18 Jun 2026, Part B Q2 & 5 \\ \hline
\end{idxtable}

\subsection*{A.3\enspace Trust, respect and harmony in the family}
\addcontentsline{toc}{section}{A.3 Trust, respect and harmony in the family}

\begin{idxtable}
Foundation of a strong human relationship &
  CIE-II, 7 May 2025, Part A Q1 & 1 \\ \hline
Disrespect Arising out of & CIE-II, 7 May 2025, Part A Q2 & 1 \\ \hline
Care is the feeling & CIE-II, 7 May 2025, Part A Q3 & 1 \\ \hline
Fulfilment of relationship means & CIE-II, 7 May 2025, Part A Q4 & 1 \\ \hline
With the complete understanding of respect &
  CIE-II, 7 May 2025, Part A Q5 & 1 \\ \hline
Relationship is between one Self (I1) and another Self (I2) &
  CIE-II, 7 May 2025, Part B Q1 & 2$+$3 \\ \hline
Distinguish between Intention and Competence with respect Trust &
  CIE-II, 7 May 2025, Part B Q2 & 5 \\ \hline
Define complementary in the context of the complete content of respect &
  CIE-II, 7 May 2025, Part B Q3 & 5 \\ \hline
Disrespect arising from differentiation leading to discrimination &
  CIE-II, 7 May 2025, Part B Q4 & 5 \\ \hline
Natural acceptance values in relationship: Reverence and Glory &
  CIE-II, 7 May 2025, Part B Q5 & 5 \\ \hline
Which of the following is a fundamental human aspiration related to
  relationships? & SEE, Jun/Jul 2025, Q1.4 & 01 \\ \hline
Foundation of a strong human relationship &
  SEE, Jun/Jul 2025, Q1.6 & 01 \\ \hline
Respect in a relationship means & SEE, Jun/Jul 2025, Q1.7 & 01 \\ \hline
Physical facility is necessary, but relationship is also necessary &
  SEE, Jun/Jul 2025, Q2b & 06 \\ \hline
Right understanding in relationships and harmony in the family &
  SEE, Jun/Jul 2025, Q3a & 07 \\ \hline
Explain the concept of harmony in the family &
  SEE, Jun/Jul 2025, Q4b & 07 \\ \hline
\end{idxtable}

\subsection*{A.4\enspace Harmony in the society --- justice, goals and values}
\addcontentsline{toc}{section}{A.4 Harmony in the society}

\begin{idxtable}
How does ensuring justice in relationships contribute to society? &
  SEE, Jun/Jul 2025, Q1.1 & 01 \\ \hline
A harmonious world is created by values at 4 levels &
  SEE, Jun/Jul 2025, Q1.9 & 01 \\ \hline
Four goals of human beings living in a society &
  SEE, Jun/Jul 2025, Q3b & 06 \\ \hline
Justice and preservation in a harmonious society &
  SEE, Jun/Jul 2025, Q4a & 06 \\ \hline
\end{idxtable}

\subsection*{A.5\enspace The self, value education and physical facility}
\addcontentsline{toc}{section}{A.5 The self, value education and physical
  facility}

\begin{idxtable}
Participation of the human being in the role of physical facility &
  Improvement CIE, Jun 2025, Quiz Q1 & 1 \\ \hline
Imagination is continuous with & SEE, Jun/Jul 2025, Q1.2 & 01 \\ \hline
Transformation from animal consciousness to human consciousness &
  SEE, Jun/Jul 2025, Q1.3 & 01 \\ \hline
Which of the following capacity leads to desires &
  SEE, Jun/Jul 2025, Q1.5 & 01 \\ \hline
Participation of the human being in the role of physical facility &
  SEE, Jun/Jul 2025, Q1.8 & 01 \\ \hline
Value of an entity, and the value of a human being &
  SEE, Jun/Jul 2025, Q2a & 08 \\ \hline
Domain of consciousness in accomplishing development or permanent change &
  Improvement CIE, 18 Jun 2026, Part B Q4 & 5 \\ \hline
\end{idxtable}

\chapter{Questions Set More Than Once}

Six topics in this collection have been examined on more than one paper. Three
were reprinted word for word; the other three were narrowed or widened, but
around the same idea. All are listed below with every occurrence and every
mark value, because a question that has already been asked twice is the most
economical thing to prepare.

\vspace{2pt}
Note what happens to the marks. \emph{Conformance in the four orders} was
worth 6 marks on a CIE and 7 on the SEE, and its Bloom's level was L4 both
times --- so the semester-end paper wanted the same analysis, one point
longer. That is the general pattern here: a repeat keeps its level and gains a
mark or two when it moves to the SEE.

\renewcommand{\arraystretch}{1.35}
\arrayrulecolor{rulegrey}
\setlength{\LTpre}{8pt}\setlength{\LTpost}{4pt}
\begin{longtable}{P{76mm} P{72mm}}
\rowcolor{slate}
\hdr{Question} & \hdr{Set in} \\ \hline \endfirsthead
\rowcolor{slate}
\hdr{Question} & \hdr{Set in} \\ \hline \endhead
\textbf{Foundation of a strong human relationship} --- identical wording and
  identical options both times &
  CIE-II, 7 May 2025, Part A Q1 (1 M, CO1, L1)\newline
  SEE, Jun/Jul 2025, Q1.6 (01 M, CO1, L2) \\ \hline
\textbf{`Conformance', and conformance in the four orders} --- identical
  wording both times &
  Improvement CIE, Jun 2025, Test Q2 (6 M, CO2, L4)\newline
  SEE, Jun/Jul 2025, Q6a (07 M, CO3, L4) \\ \hline
\textbf{Participation of the human being in the role of physical facility}
  --- reworded: \emph{`in nurture, protection and providing means for the
  body'} became \emph{`to help and preserve its utility'}; the four options
  are unchanged &
  Improvement CIE, Jun 2025, Quiz Q1 (1 M, CO1, L2)\newline
  SEE, Jun/Jul 2025, Q1.8 (01 M, CO3, L2) \\ \hline
\textbf{The four orders of nature} --- asked three times, at three different
  lengths: as a five-mark explanation, as a six-mark question with a schematic
  diagram, and as a one-mark list &
  Improvement CIE, 28 Aug 2024, Part B Q3 (5 M, CO4, L2)\newline
  Improvement CIE, Jun 2025, Test Q1 (6 M, CO1, L2)\newline
  Improvement CIE, 18 Jun 2026, Part A Q2 (1 M, CO2, L1) \\ \hline
\textbf{Svabhava --- the natural characteristic} --- narrowed and widened
  across three papers: the human order alone, then all four orders, then the
  self &
  Improvement CIE, 28 Aug 2024, Part B Q2 (5 M, CO2, L1)\newline
  SEE, Jun/Jul 2025, Q5b (07 M, CO4, L3)\newline
  Improvement CIE, 18 Jun 2026, Part A Q3 (1 M, CO2, L2) \\ \hline
\textbf{Units and space} --- asked as a seven-mark elaboration and then as a
  one-mark differentiation &
  Improvement CIE, Jun 2025, Test Q4 (7 M, CO4, L4)\newline
  Improvement CIE, 18 Jun 2026, Part A Q5 (1 M, CO4, L2) \\ \hline
\end{longtable}

\begin{note}
\textbf{\sffamily What the repeats tell you.} Six topics carry every repeat in
this collection: the four orders, svabhava, conformance, units and space, the
utility value of physical facility, and trust as the foundation of
relationship. Between them they account for 15 of the 59 questions in this
book. If revision time is short, that is where it goes first.
\end{note}

\end{document}
